Four findings, in descending order of how well they are supported.

\paragraph{1. Under the published cost model, both arms substantially exceed the
published result.} Arm B reaches an annualised excess return of $+12.75\%$ with
an information ratio of $1.71$ and a maximum drawdown of $-9.01\%$, against the
published $+4.68\%$, $0.65$ and $-11.80\%$; Arm A reaches $+7.91\%$ and $1.15$.
Both clear the four hand-written base features and Qlib's own factor subset on
every axis. Arm B beats Arm A in all four calendar years, and turns 2025 --- the
year in which the base features lose $12.38\%$ --- into a flat $-0.21\%$. These
differences are far outside combiner-seed noise, resolvable at 17--20$\times$
on return and over 100$\times$ on IC.

\paragraph{2. The advantage of the net-of-cost objective is consistent, and its
size is roughly five points of annualised return.} The paired difference
reproduced across two independent runs at $+5.23$pp and $+4.84$pp on return, and
$+0.0120$ and $+0.0123$ on IC. That consistency is more informative than either
run alone, because the individual arms moved by four points and $0.013$
respectively between the same two runs. The pairing removes a common-mode
run-to-run shift that is an order of magnitude larger than anything the seed
error bars capture (\S\ref{sec:pairing}).

\paragraph{3. The utility's selection mechanism contributed nothing detectable.}
This is the sharpest negative result here, and it was only measurable because the
first run's utility was inert. Repairing it changed the mechanism unambiguously
--- the repository now accumulates across process boundaries and $\util$ takes
distinct values per batch instead of a constant $0.9500$ --- yet the paired
effect was essentially unchanged: $+5.23$pp before the fix, $+4.84$pp after.
Whatever advantage the treatment has, it was already fully present when the
objective was ranking nothing at all.

The mechanism that remains is the feedback channel. Net-of-cost metrics reach the
generator as text regardless of whether $\util$ ranks anything, and the arms
diverge in exactly the way that channel would predict: Arm A mines one- to
five-day microstructure, Arm B mines three- to twelve-month horizons with
explicit liquidity and crowding terms. On the evidence here, telling the model
about costs is doing the work; ranking candidates by a cost-aware utility is not
adding to it.

\paragraph{4. Under a realistic cost model, none of this is tradeable.} Every
configuration loses money: $-6.51\%$ (Arm A), $-7.12\%$ (Arm B) and $-6.92\%$
(no mining at all), at roughly $4.2$ basis points per day, or about $10\%$
annually. The arms are indistinguishable from each other and from the baseline
--- the Arm B minus Arm A difference is $1.0\times$ the seed noise --- and both
mined libraries score \emph{below} the no-mining baseline on Rank IC. The
$+12.75\%$ headline and the $-7.12\%$ net figure are the same factors under two
different assumptions about friction, and the distance between them is the whole
of the slippage and impact the published methodology does not charge.

Two structural reasons explain why the capacity-aware objective cannot rescue
this. Book turnover is pinned near $\NDrop/\TopK$ by the portfolio construction,
so both arms trade the same amount and the objective has almost no room to
express itself by trading less; measured cost differs between them by $0.04$
basis points per day. And the mined libraries carry an in-sample to
out-of-sample gap around $0.15$, against $0.03$ for the hand-written baseline ---
roughly five times the overfitting, which is what searching a large expression
space against a fixed evaluation buys.

\subsection*{What would settle the open questions}

The effect estimate is well behaved but rests on two pairs. Replicating across
generation seeds --- \texttt{QA\_SEED} varies the trajectory, and the driver
already runs both arms from one invocation --- would convert a consistent point
estimate into a tested one. Three to five paired invocations would give a
standard deviation on the paired difference; the measurements here suggest that
difference is stable enough for a small number of replicates to be informative.

Separating the feedback channel from the selection channel deserves a direct
test rather than an accident. An arm that receives net-of-cost feedback but
selects by RankIC would isolate the contribution that finding~3 attributes to
feedback, and it costs one additional arm rather than a redesign.

Finally, the tradeability result is a statement about this portfolio
construction as much as about these factors. Turnover is structurally capped, so
a construction with a genuine turnover budget --- the mean--variance member of
the admissible family, with an explicit cap --- is the only way to discover
whether a capacity-aware objective can produce a book that survives its own
trading costs. Until then, the honest summary is that the net-of-cost objective
produces measurably better signal under the published accounting, and that no
configuration tested here is tradeable under a realistic one.
